\section{Conclusion}
\label{sec:ieee-conclusion}

ORIUS argues that the right universal object for Physical AI safety is not a single
planner, plant, or benchmark. It is a runtime safety layer that mediates the gap
between degraded observation and physical actuation. Once observation quality becomes a
source of safety failure, the system must expose reliability, uncertainty inflation,
action repair, fallback, and certificate semantics as first-class runtime objects.

That is the role of the Detect--Calibrate--Constrain--Shield--Certify kernel. It turns
OASG from an intuitive failure mode into a measurable and governable runtime event, and
it does so under one contract that can be instantiated across batteries, autonomous
vehicles, industrial plants, healthcare monitoring, navigation, and aerospace without
pretending those domains share the same nominal controller or safety geometry.

The evidence claim is intentionally disciplined. Battery is the witness row. AV,
industrial, and healthcare are defended bounded rows. Navigation and aerospace remain
the two visible blockers for equal-domain closure. That asymmetry does not weaken the
architectural claim. It is what makes the paper scientifically credible: ORIUS can be
universal in runtime contract while still governed by a public promotion gate.

The field-level implication is larger than any single benchmark number. ORIUS proposes
that Physical AI should be organized around a reusable runtime safety layer that sits
between nominal intelligence and actuation, and that this layer should be audited
through one benchmark schema and one governance surface. If the remaining gated rows are
closed, the same architecture can support a stronger universality claim. Even before
that point, ORIUS establishes a reusable safety-layer grammar for Physical AI under
degraded observation.
